\def\topic{Schr\"odinger semigroups with bounded potentials}
\ifdefined\MyHeaderIncluded
\else
\def\qwertyflag{true}

\ifdefined\MyHeaderIncluded

\else
  \newcommand{\foobar}{FooBar}
  \def\MyHeaderIncluded{true}
% if is till the end of the file

\documentclass[10pt,a4paper]{article}
\usepackage[utf8]{inputenc}
%\usepackage[german]{babel}
%\usepackage[cyrillic]{babel}
\usepackage[T1]{fontenc}
%\usepackage[T2A]{fontenc}
%\usepackage[english, russian]{babel}
\usepackage{amsmath}
\usepackage{amsfonts}
\usepackage{amssymb}
\usepackage{amsthm}
\usepackage{mathtools} %coloneqq

\usepackage{bussproofs} %

%\renewcommand{\thesubsection}{\arabic{subsection}}
%\newtheorem{thm}{Theorem}[section]
%\renewcommand{\thethm}{\arabic{subsection}.\arabic{thm}}

\usepackage{thmtools}
\usepackage{hyperref}
\hypersetup{
    colorlinks=true,
    linkcolor=blue,
    filecolor=magenta,      
    urlcolor=cyan,
    pdftitle={Overleaf Example},
    %pdfpagemode=FullScreen,
    }
%\usepackage{stix-type1}%\turnediota
\theoremstyle{definition}

\newtheorem{constructor}{Constructor}[section]
\newtheorem{iaxiom}{Axiom}
\newtheorem{axiom}{Axiom}
\newtheorem*{axiom*}{Axiom}
\newtheorem{definition}{Definition}[section]
\newtheorem*{definition*}{Definition}
\newtheorem{remark}{Remark}[section]
\newtheorem*{remark*}{Remark}
\newtheorem{recall}{Recall}[section]
\newtheorem*{recall*}{Recall}

\newtheorem{lemma}{Lemma}[section]
\newtheorem{theorem}{Theorem}[section]
\newtheorem*{theorem*}{Theorem}
\newtheorem{note}{Note}[section]
\usepackage{etoolbox}

%\newcommand*{\@Dashv}[2]{%
%  \reflectbox{$\m@th#1#2$}%
%}
\newcommand*{\separation}[3]{% 1)binder, 2)set, 3)condition
  \{#1\in #2: #3\}
}

% https://www.freecodecamp.org/news/simple-regex-tricks-for-beginners-3acb3fa257cb/
%  \\mathcal\{T\} |---> \topology
%  \\mathcal\{([^\}]*)\}   |--->  $1
\newcommand{\AND}{\ \ \&\ \ }
\newcommand{\any}{\mathrm{\ any}}
\newcommand{\borel}{\mathcal{B}}
\newcommand{\can}{\mathrm{can}}
\newcommand{\canonize}{\mathrm{canonize}}
\newcommand{\closure}{\mbox{cl}}
\newcommand{\complex}{\mathbb{C}}
\newcommand{\converges}{\longrightarrow}
\newcommand{\countable}{\mbox{\ countable}}
\newcommand{\ldefined}{\mbox{\ defined}}
\newcommand{\dom}{\mbox{dom}}
\newcommand{\Dom}{\mbox{Dom}}
\newcommand{\Done}{\noindent \textbf{Done.}\\}
\newcommand{\lexists}{\mbox{\ exists}}
\newcommand{\lforall}{\mbox{\ forall}}
\newcommand{\lunique}{\mbox{\ unique}}
\newcommand{\family}{\mathrm{\ family}}
\newcommand{\Fm}{\mathrm{Fm}}
\newcommand{\FuncSymb}{\mathrm{FuncSymb}}
\newcommand{\function}{\mathrm{\ function}}
\newcommand{\functional}{\mbox{\ functional}}
\newcommand{\IFF}{\ \Longleftrightarrow\ }
\renewcommand{\iff}{\ \longleftrightarrow\ }
\newcommand{\IMPLIES}{\ \Longrightarrow\ }
\renewcommand{\implies}{\ \longrightarrow\ }
\newcommand{\Inf}{\mathrm{Inf}}
\newcommand{\injective}{\mbox{\ injective}}
\newcommand{\injection}{\mbox{\ injection}}
\newcommand{\interior}{\mbox{int}}
\newcommand{\intersection}{\bigcap}
\newcommand{\isempty}{\mathrm{\ empty}}
\newcommand{\LB}{\mathrm{LB}}
\newcommand{\nghbset}{\mathcal{N}}
\newcommand{\nonempty}{\mathrm{\ nonempty}}
\newcommand{\Ord}{\mathrm{Ord}}
\newcommand{\power}{\mathcal{P}}
\newcommand{\pr}{\mathrm{pr}}
\newcommand{\fst}{\mathrm{fst}}
\newcommand{\snd}{\mathrm{snd}}
\newcommand{\PredSymb}{\mathrm{PredSymb}}
\newcommand{\Proof}{\noindent \textbf{Proof.}\\}
\newcommand{\ran}{\mathrm{ran}}
\newcommand{\Ran}{\mathrm{Ran}}
\newcommand{\real}{\mathbb{R}}
\newcommand{\relation}{\mathrm{\ relation}}
\newcommand{\set}{\mathrm{\ set}}
\newcommand{\singleton}{\mathrm{\ singleton}}
\newcommand{\semantics}{\mathrm{semantics}}
\newcommand{\semialgebra}{\mathcal{S}}
\newcommand{\Sup}{\mathrm{Sup}}
\newcommand{\surjective}{\mbox{\ surjective\ on\ }}
\newcommand{\Term}{\mathrm{Term}}
\newcommand{\Type}{\mathrm{Type}}
\newcommand{\topology}{\mathcal{T}}
\newcommand{\total}{\mbox{\ total\ on\ }}
\newcommand{\transitive}{\mathrm{\ transitive}}
\newcommand{\union}{\bigcup}
\newcommand{\lTrue}{\mathrm{True}}
\newcommand{\lFalse}{\mathrm{False}}
\newcommand{\lUncertain}{\mathrm{Uncertain}}
\newcommand{\lif}{\mathrm{if}}
\newcommand{\lthen}{\mathrm{then}}
\newcommand{\lelse}{\mathrm{else}}
\newcommand{\lexception}{\mathrm{exception}}
\newcommand{\lundef}{\mathrm{undef}}
\newcommand{\univocal}{\mathrm{\ univocal}}
\newcommand{\univocalat}{\mathrm{\ univocal\ at\ }}
\newcommand{\UB}{\mathrm{UB}}
\newcommand{\Un}{\mathcal{U}}
\newcommand{\Var}{\mathrm{Var}}

\newcommand{\notvarphi}{\phi}
\renewcommand{\phi}{\varphi}

\newcommand{\comment}[1]{\quad\mbox{(#1)}}
% local commands
%\preto\section{\undef\blipo} %etoolbox
\preto\subsection{\undef\blipo} %etoolbox
% https://tex.stackexchange.com/questions/105684/limit-macro-newcommand-scope-to-section
\usepackage{graphicx}
\usepackage{authblk}

\setlength\parindent{0pt} % whole doc of \noindent
\usepackage{soul}
%\usepackage{mathgl}

\fi % guard
\ifdefined\topic
\else
\def\topic{No Titel}
\fi
\title{\topic}
\author[1]{Georgy Dunaev}
\begin{document}
\maketitle
\fi
\section{Schr\"odinger semigroups with bounded potentials}
\label{sec:bounded_schrodinger}
This chapter works on $\mathbb R^d$, for a positive integer $d$, with Lebesgue measure
and complex $L^2$. Its analytic inputs are the Fourier--Plancherel
theorem, the spectral calculus for self-adjoint operators, and the
usual product-integration theorems. We take these standard analytic
results as prerequisites and prove the bounded-potential perturbation
and kernel statements below directly. In particular, no pointwise limit of product formulas
is used to assert existence or continuity of a kernel.
We use the inner product $\langle f,g\rangle=\int f\overline g$,
linear in its first variable. An operator norm written without
subscripts means the norm from $L^2(\mathbb R^d)$ to itself.
The imported spectral calculus says that a self-adjoint operator $A$
has a projection-valued spectral measure $E_A$ with
$A=\int\lambda\,dE_A(\lambda)$ on the domain determined by
$\int\lambda^2\,d\langle E_A(\lambda)f,f\rangle<\infty$;
bounded Borel functions of $A$ act by the corresponding spectral
integrals. Fourier--Plancherel says that the normalized transform below
extends to a unitary operator on $L^2$. These standard background
theorems are used, not reproved here; see \parencite{Weidmann1980Operators}.
The concrete joint-continuity conclusion is also treated in
\parencite[Theorem~B.7.1]{Simon1982Semigroups}; the bounded-potential proof
below supplies its own argument, distinct from the abstract
separate-continuity criterion in Simon's Section~A.1.

\subsection{The operator and its sign convention}
Use the unitary Fourier transform
\[
\widehat f(\xi)=(2\pi)^{-d/2}\int_{\mathbb R^d}
                   e^{-ix\cdot\xi}f(x)\,dx
\]
initially on integrable square-integrable functions and then by
Plancherel extension. Define the nonnegative Laplacian by
\[
H_0=-\sum_{j=1}^d\partial_j^2,\qquad
D(H_0)=\{f\in L^2:|\xi|^2\widehat f(\xi)\in L^2\},\qquad
\widehat{H_0f}=|\xi|^2\widehat f.
\]
The differential expression agrees with this definition on smooth
functions of compact support. Throughout, $V$ is a real essentially
bounded measurable function, $M=\|V\|_\infty$, and multiplication by
$V$ is denoted by the same symbol. Choose a representative with
$|V(x)|\leq M$ everywhere. None of the resulting operators or kernels
depends on this choice on null sets.

\begin{proposition}[Bounded real perturbations]
\label{prop:bounded_schrodinger_selfadjoint}
The operator $H=H_0+V$ on $D(H_0)$ is self-adjoint and bounded below
by $-M$. Its spectral semigroup $T_t=e^{-tH}$, $t\geq0$, is strongly
continuous and satisfies $\|T_t\|_{2\to2}\leq e^{Mt}$.
\end{proposition}
\Proof
The real multiplication operator $u\mapsto|\xi|^2u$ on its maximal
domain is self-adjoint. To check the domain assertion, if $v$ belongs
to its adjoint domain with adjoint value $w$, test against all $L^2$
functions supported in $\{|\xi|\leq n\}$. On that set the multiplier
is bounded, so these tests give $w=|\xi|^2v$ there. Letting $n$ increase
proves $|\xi|^2v\in L^2$, exactly the original domain. The reverse
inclusion follows by the integral pairing. Plancherel conjugation
therefore makes $H_0$ self-adjoint and nonnegative; its domain is dense
by frequency truncation.

Multiplication by real $V$ is bounded self-adjoint. In general, for a
densely defined operator $A$ and bounded operator $B$, a vector belongs
to $D((A+B)^*)$ exactly when it belongs to $D(A^*)$: the bounded
functional contributed by $B$ can be subtracted from the defining
adjoint identity. Thus $(A+B)^*=A^*+B^*$ on that domain. Applying
this with $A=H_0$ proves self-adjointness of $H$. Also
\[
\langle Hf,f\rangle
 =\int |\xi|^2|\widehat f(\xi)|^2\,d\xi
   +\int V(x)|f(x)|^2\,dx\geq-M\|f\|_2^2.
\]
The spectral theorem now gives $T_t$, its semigroup law and norm bound.
Strong continuity follows by dominated convergence in each vector's
spectral measure, since $e^{-t\lambda}\to1$ and the spectrum lies in
$[-M,\infty)$.
\Done

This is the bounded-perturbation case of the self-adjoint perturbation
principle: $\|Vf\|_2\leq M\|f\|_2$ gives relative bound zero with
respect to $H_0$. The direct adjoint-domain proof avoids needing the
more general Kato--Rellich theorem for unbounded relatively bounded
perturbations; see \parencite{Kato1995Perturbation} for that broader setting.

Put $P_t=e^{-tH_0}$. The Gaussian Fourier calculation gives
\[
(P_tf)(x)=\int_{\mathbb R^d}p_t(x-y)f(y)\,dy,\qquad
p_t(r)=(4\pi t)^{-d/2}\exp\bigl(-|r|^2/(4t)\bigr),\quad t>0.
\]
Indeed the Fourier multiplier is $e^{-t|\xi|^2}$;
Lemma~\ref{lem:gaussian_fourier_multiplier} computes the Gaussian
transform with precisely this normalization. In particular
\[
p_t\geq0,\qquad \int p_t=1,\qquad
\int p_{t-s}(x-z)p_s(z-y)\,dz=p_t(x-y)\quad(0<s<t).
\]
The last identity follows either by the Fourier multipliers or by
the square-completion calculation in Lemma~\ref{lem:free_gaussian_convolution}.
Notice the sign:
the heat generator is $-H_0=\sum_j\partial_j^2$.

\subsection{Duhamel's equation and its kernel series}
All time integrals below are taken after applying the operator to a
fixed vector. For a continuous map $u:[a,b]\to L^2$, its integral is
the norm limit of Riemann sums: uniform continuity makes these sums
Cauchy, and completeness of $L^2$ supplies their limit. The triangle
inequality for sums passes to the limit and gives
\[
\left\|\int_a^b u(s)\,ds\right\|_2
 \leq\int_a^b\|u(s)\|_2\,ds.
\]
Bounded linear maps commute with this integral by passage through the
sums. The same inequality and scalar dominated convergence justify
the limits of continuous vector-valued integrands used below; the
fundamental theorem for a continuously differentiable $L^2$-valued
map follows by the usual Riemann-sum argument. Thus no operator-norm
continuity of the free semigroup at zero is required to form these
strong integrals.

\begin{lemma}[The bounded-perturbation expansion]
\label{lem:schrodinger_duhamel}
For $t\geq0$,
\[
T_t=P_t-\int_0^t P_{t-s}VT_s\,ds
     =\sum_{n=0}^{\infty}A_n(t),
\]
where $A_0(t)=P_t$ and
\[
A_n(t)=-\int_0^t P_{t-s}V A_{n-1}(s)\,ds,\qquad
\|A_n(t)\|_{2\to2}\leq\frac{M^nt^n}{n!}.
\]
The integrals are strong integrals after application to an $L^2$
vector, and the series converges in operator norm, uniformly on bounded
time intervals.
\end{lemma}
\Proof
For $f\in D(H_0)=D(H)$, differentiate $P_{t-s}T_sf$ for $0<s<t$.
The spectral calculus preserves $D(H)$ under $T_s$, and
$H_0T_sf=HT_sf-VT_sf$ is continuous in $s$ in $L^2$. The product rule
therefore gives
\[
\frac{d}{ds}(P_{t-s}T_sf)=-P_{t-s}VT_sf.
\]
Integrating and using strong continuity at the endpoints proves the
first equation on the dense domain. The uniform operator bound on the
integrand extends it to every $L^2$ vector.

The recursive integrals defining $A_n$ exist inductively by strong
continuity. More explicitly, substitute $s=t\theta$ on the fixed
interval $0\leq\theta\leq1$. The maps
$P_{t(1-\theta)}VA_{n-1}(t\theta)f$ are continuous in $t$ for each
$\theta$, by the induction hypothesis and the uniform bound on $P$.
On bounded time intervals the induction bound supplies an integrable
constant majorant. Dominated convergence therefore gives strong
continuity of $A_n(t)$; its value at zero is zero for $n\geq1$.
Induction and $\|P_s\|\leq1$ give
\[
\|A_n(t)\|\leq M\int_0^t\frac{M^{n-1}s^{n-1}}{(n-1)!}\,ds
 =\frac{M^nt^n}{n!}.
\]
Consequently the series converges with the stated uniformity and
satisfies the same integral equation, by termwise integration.
For uniqueness, the difference $D_t$ of two locally bounded operator
solutions satisfies $D_t=-\int_0^tP_{t-s}VD_s\,ds$.
On $0\leq t\leq b$, iteration bounds its norm by
$C_b M^nt^n/n!$ for every $n$, where $C_b=\sup_{s\leq b}\|D_s\|$.
This tends to zero. Thus the sum equals $T_t$.
\Done

Define measurable kernels, for $t>0$, by
\[
k_0(t,x,y)=p_t(x-y),\qquad
k_n(t,x,y)=-\int_0^t\int_{\mathbb R^d}
 p_{t-s}(x-z)V(z)k_{n-1}(s,z,y)\,dz\,ds.
\]
Values at the two time endpoints need not be assigned in these
Lebesgue integrals. Induction using Gaussian convolution proves absolute
convergence and the pointwise estimate
\[
|k_n(t,x,y)|\leq\frac{M^nt^n}{n!}p_t(x-y).
\tag{K$_n$}\label{eq:schrodinger_dyson_bound}
\]
For the induction step, the spatial integral after taking absolute
values is bounded by
\[
\frac{M^n s^{n-1}}{(n-1)!}p_t(x-y).
\]
Integration over $s$ gives the result.
Changing $V$ on a null set changes none of these integrals.

\begin{theorem}[A continuous Gaussian-bounded kernel]
\label{thm:bounded_schrodinger_kernel}
For real $V\in L^\infty(\mathbb R^d)$ the function
\[
k_V(t,x,y)=\sum_{n=0}^{\infty}k_n(t,x,y)
\]
is jointly continuous on $(0,\infty)\times\mathbb R^d\times\mathbb R^d$
and is an integral kernel for $e^{-tH}$, with
\[
|k_V(t,x,y)|\leq e^{Mt}p_t(x-y).
\]
For every $f\in L^2$, the integral
$\int k_V(t,x,y)f(y)\,dy$ is absolutely convergent at every $x$ and
gives a continuous representative of $T_tf$.
\end{theorem}
\Proof
We first prove continuity of each term; measurability of $V$ alone is
sufficient. The assertion is clear for $k_0$. In the recursion for
$k_n$, put $s=t\theta$, and restrict first to
$\varepsilon\leq\theta\leq1-\varepsilon$. On a compact set of
parameters with $0<a\leq t\leq b$ and bounded $x,y$, both Gaussian
time parameters are bounded away from zero. The induction hypothesis
makes the integrand continuous in $t,x,y$ for each $\theta,z$, even
though its fixed factor $V(z)$ is only measurable. Estimate
\eqref{eq:schrodinger_dyson_bound} bounds it by an integrable Gaussian
envelope $C\exp(-c|z|^2)$, with constants depending on this compact
set and $\varepsilon$. Dominated convergence proves continuity of
the truncated integral.

For the two omitted time intervals, the convolution estimate gives
the uniform bound
\[
\frac{M^nt^n}{(n-1)!}p_t(x-y)
 \int_{(0,\varepsilon)\cup(1-\varepsilon,1)}
                      \theta^{n-1}\,d\theta.
\]
On the stated compact parameter set this tends to zero with
$\varepsilon$. Thus $k_n$ is a locally uniform limit of continuous
functions and is continuous. Finally, on $a\leq t\leq b$,
\[
|k_n(t,x,y)|\leq\frac{(Mb)^n}{n!}(4\pi a)^{-d/2}.
\]
The series is therefore locally uniformly convergent and has a
jointly continuous sum. Summing the sharper Gaussian estimates gives
the asserted bound.

We will use the following elementary density fact. For $1\leq p<\infty$
and $f\in L^p(\mathbb R^d)$, put
\[
f_m=f\mathbf1_{\{|x|\leq m,\ |f(x)|\leq m\}}.
\]
Each $f_m$ is bounded and compactly supported. The functions
$|f-f_m|^p$ tend to zero almost everywhere and are bounded by $|f|^p$,
so dominated convergence gives $\|f-f_m\|_p\to0$.
In particular these functions are dense in both $L^1$ and $L^2$,
and $L^1\cap L^2$ is dense in $L^1$. This argument does not claim
density in $L^\infty$.

To identify the operator, first take a bounded compactly supported
function $f$. The same estimates and Tonelli's theorem for absolute
values justify interchanging all integrals in the recursion; thus
$A_n(t)f(x)=\int k_n(t,x,y)f(y)\,dy$ almost everywhere, inductively.
They also justify summation under this integral. The integral operator
defined by the sum is bounded on $L^2$: its absolute value is at most
$e^{Mt}P_t|f|$, whose $L^2$ norm is at most $e^{Mt}\|f\|_2$.
Density and the operator-norm series in Lemma~\ref{lem:schrodinger_duhamel}
identify it with $T_t$ on all of $L^2$.

For an arbitrary $f\in L^2$, Cauchy--Schwarz and the Gaussian bound
give absolute convergence at each $x$. On any bounded set of $x$,
the rows $k_V(t,x,\cdot)$ have a common Gaussian envelope in $L^2$.
Their pointwise continuity and dominated convergence therefore imply
continuity in $L^2$ as $x$ varies. Pairing these rows with $f$ proves
continuity of the displayed representative. Its pointwise bound also
gives $\|T_tf\|_\infty\leq e^{Mt}\|p_t\|_2\|f\|_2$.
\Done

\paragraph{The integral inequality used for smoothing.}
For clarity, $L^p=L^p(\mathbb R^d)$ consists of measurable functions
modulo equality almost everywhere, with
\[
\|f\|_p=\left(\int|f|^p\right)^{1/p}\quad(1\leq p<\infty),
\qquad \|f\|_\infty=\operatorname*{ess\,sup}|f|.
\]
All operator bounds below refer to these norms. The kernel supplies
specified continuous representatives when its spatial regularity is used.

For conjugate exponents $1<\alpha,\beta<\infty$, meaning
$1/\alpha+1/\beta=1$, elementary
one-variable minimization gives
\[
ab\leq a^\alpha/\alpha+b^\beta/\beta\qquad(a,b\geq0).
\]
Indeed the minimum of $a^\alpha/\alpha-ab$ as a function of $a$
is $-b^\beta/\beta$. Apply this inequality to
$|u|/\|u\|_\alpha$ and $|v|/\|v\|_\beta$ and integrate to obtain
H\"older's inequality
$\int|uv|\leq\|u\|_\alpha\|v\|_\beta$.
Zero norms are handled directly; at the conjugate endpoints $1,\infty$
the inequality follows from the almost-everywhere essential bound.

We also use its three-factor form: if $r_j\in[1,\infty]$ and
$\sum_{j=1}^3 1/r_j=1$, then
$\int|u_1u_2u_3|\leq\prod_{j=1}^3\|u_j\|_{r_j}$.
For completeness, put $1/r=1/r_1+1/r_2$. When $r<\infty$,
apply the two-factor inequality to $|u_1|^r,|u_2|^r$ with conjugate
exponents $r_1/r,r_2/r$, giving
$\|u_1u_2\|_r\leq\|u_1\|_{r_1}\|u_2\|_{r_2}$.
For $r=\infty$ this follows directly from essential bounds.
A second application with conjugate exponents $r,r_3$ proves the
claim, including all endpoint cases.

\begin{corollary}[Quantitative smoothing between Lebesgue spaces]
\label{cor:schrodinger_Lp_Lq}
Let $1\leq p\leq q\leq\infty$, and define $r\in[1,\infty]$ by
$1+1/q=1/p+1/r$, with $1/\infty=0$. The kernel defines a consistent
bounded operator $T_t:L^p\to L^q$ extending the spectral operator on
$L^p\cap L^2$, and
\[
\|T_tf\|_q\leq
c_r e^{Mt}(4\pi t)^{-\frac d2(\frac1p-\frac1q)}\|f\|_p,
\qquad
c_r=\begin{cases}r^{-d/(2r)},&r<\infty,\\1,&r=\infty.\end{cases}
\]
In particular the estimate includes $p=1,q=\infty$ and
$p=q=\infty$. No strong-continuity assertion on $L^\infty$ at time
zero is needed here. This is the bounded-potential instance of the
smoothing property in \parencite[Theorem~2.1]{CyconFroeseKirschSimon1987}.
\end{corollary}
\Proof
The Gaussian bound gives $|T_tf|\leq e^{Mt}p_t*|f|$. We record the
convolution estimate rather than requiring an interpolation theorem.
For finite $q$, the exponent relation gives $p,r\leq q$. H\"older's
inequality, applied to the three factors
\[
\bigl(|f(y)|^{p/q}|g(x-y)|^{r/q}\bigr),\quad
|f(y)|^{1-p/q},\quad |g(x-y)|^{1-r/q},
\]
with reciprocal exponents $1/q$, $1/p-1/q$, and $1/r-1/q$, gives
\[
|(g*f)(x)|^q
\leq\left(\int |f(y)|^p|g(x-y)|^r\,dy\right)
       \|f\|_p^{q-p}\|g\|_r^{q-r}.
\]
Factors whose reciprocal exponent is zero are omitted. Integrating in
$x$ and using Tonelli gives $\|g*f\|_q\leq\|g\|_r\|f\|_p$.
Zero functions are immediate; truncation justifies the calculation if
an intermediate pointwise integral is initially infinite. For
$q=\infty$, the relation says $1/p+1/r=1$, and the same estimate
is H\"older directly, including its $L^1$--$L^\infty$ endpoints.

Gaussian integration gives
\[
\|p_t\|_r=r^{-d/(2r)}(4\pi t)^{-\frac d2(1-1/r)}
\quad(r<\infty),\qquad
\|p_t\|_\infty=(4\pi t)^{-d/2}.
\]
This proves the displayed bound and absolute convergence almost
everywhere of the defining kernel integral for every $f\in L^p$.
The kernel is the same for all these spaces, so the operators agree
wherever their domains overlap. On $L^p\cap L^2$, the preceding
theorem identifies that integral with the spectral semigroup.
\Done

\subsection{Pointwise identities and positivity}
An integral operator has a unique continuous kernel within the
Gaussian-bounded class just constructed. Indeed, equality of two such
operators gives zero integral of the kernel difference against every
product of bounded compactly supported test functions in $x$ and $y$.
If a continuous difference were nonzero at a point, multiplication by
a scalar of modulus one and sufficiently small product neighbourhoods
would make its real part have one strict sign there. Their indicator
tests give a contradiction. This observation upgrades operator
identities to pointwise kernel identities; it is stronger than merely
choosing unrelated almost-everywhere representatives at each time.

\begin{proposition}[Symmetry, composition, and two-sided bounds]
\label{prop:schrodinger_kernel_identities}
For $s,t>0$ and $x,y\in\mathbb R^d$,
\[
\begin{gathered}
k_V(t,x,y)=k_V(t,y,x)\in\mathbb R,\\
k_V(t+s,x,y)=\int_{\mathbb R^d}
                     k_V(t,x,z)k_V(s,z,y)\,dz,\\
e^{-Mt}p_t(x-y)\leq k_V(t,x,y)\leq e^{Mt}p_t(x-y).
\end{gathered}
\]
In particular the semigroup preserves positivity.
\end{proposition}
\Proof
Every term in the kernel series is real. Self-adjointness of $T_t$
identifies its kernel with the transposed real kernel, and continuous
kernel uniqueness gives symmetry everywhere. The spatial composition
integral converges absolutely by Gaussian convolution. It is continuous
in $x,y$: on compact parameter sets the product has an integrable
Gaussian envelope in $z$, so dominated convergence applies. Fubini
identifies its operator with $T_tT_s=T_{t+s}$. Kernel uniqueness gives
the second equation pointwise.

For positivity put $W=M-V$, so $0\leq W\leq2M$ almost everywhere.
Apply the preceding kernel construction to the potential $-W$.
All of its Dyson terms are nonnegative, since each recursion has
factor $-(-W)=W\geq0$. Its kernel is therefore at least $p_t$.
The operators satisfy
\[
H_0+V=H_0-W+M,\qquad
e^{-t(H_0+V)}=e^{-Mt}e^{-t(H_0-W)}.
\]
Continuous kernel uniqueness gives $k_V=e^{-Mt}k_{-W}$, and hence
the lower bound. The upper bound is the already proved absolute
Gaussian bound. Integrating the nonnegative kernel against any
nonnegative $L^2$ function proves positivity preservation.
\Done

\paragraph{A normalization check: constant potentials.}
If $V(x)=c$ for a real constant $c$, spectral calculus gives
$e^{-t(H_0+c)}=e^{-ct}P_t$. Continuous kernel uniqueness therefore gives
\[
k_c(t,x,y)=e^{-ct}(4\pi t)^{-d/2}
                     \exp\bigl(-|x-y|^2/(4t)\bigr).
\]
The same formula follows directly from the series: Gaussian convolution
gives $k_n=(-c)^nt^np_t(x-y)/n!$, and summation gives $e^{-ct}p_t$.
For a prescribed $M\geq0$, the choices $c=M$ and $c=-M$ attain,
respectively, the lower and upper Gaussian bounds. Thus the time factor
and the sign of the potential in those bounds cannot be improved
uniformly over all real potentials with $\|V\|_\infty\leq M$.

\subsection{The heat equation in its precise solution class}
For a merely measurable potential, a continuous kernel need not give
classical second spatial derivatives. The natural equation here is an
$L^2$ evolution equation, whose spatial differential expression is
understood in distributions. As in the free heat chapter, write
$L=\sum_{j=1}^d\partial_j^2$, so that $H_0=-L$.
If $f\in D(H_0)$ and $\phi\in C_c^\infty(\mathbb R^d)$, integration
by parts in the Fourier transform gives
$\widehat{-L\phi}=|\xi|^2\widehat\phi$. Plancherel then gives
\[
\langle H_0f,\phi\rangle
=\langle |\xi|^2\widehat f,\widehat\phi\rangle
=\langle\widehat f,|\xi|^2\widehat\phi\rangle
=\langle f,-L\phi\rangle.
\]
All pairings are finite by Cauchy--Schwarz: both $H_0f$ and
$-L\phi$ belong to $L^2$. This identity is exactly the
distributional interpretation of $H_0f=-L f$, without requiring
pointwise derivatives of $f$.

\begin{proposition}[The strong solution and its distributional equation]
\label{prop:schrodinger_heat_solution}
For every $f\in L^2$ and $b>0$, $u(t)=T_tf$ is the unique function in
\[
C([0,b];L^2)\cap C^1((0,b);L^2)
\]
such that $u(0)=f$, $u(t)\in D(H)=D(H_0)$ for $0<t<b$, and
$u'(t)=-Hu(t)$ in $L^2$. At every positive time,
\[
\|HT_t\|_{2\to2}\leq Me^{Mt}+\frac1{et}.
\]
Its kernel representative is jointly continuous in $(t,x)$ for $t>0$.
For every $\phi\in C_c^\infty(\mathbb R^d)$ it satisfies
\[
\frac{d}{dt}\int u(t,x)\overline{\phi(x)}\,dx
=\int u(t,x)\overline{L\phi(x)}\,dx
 -\int V(x)u(t,x)\overline{\phi(x)}\,dx.
\]
The initial condition is convergence to $f$ in $L^2$.
\end{proposition}
\Proof
The spectrum of $H$ is contained in $[-M,\infty)$. For $\lambda\geq0$,
$\lambda e^{-t\lambda}\leq1/(et)$, whereas for
$-M\leq\lambda<0$ one has
$|\lambda|e^{-t\lambda}\leq Me^{Mt}$. Thus spectral calculus gives
the displayed bound. It also gives
\[
\int\lambda^2e^{-2t\lambda}\,
             d\langle E_H(\lambda)f,f\rangle<\infty,
\]
so $T_tf\in D(H)$ for every $t>0$.
On any compact positive-time interval the scalar multipliers
$\lambda e^{-t\lambda}$ have a common bound. The difference quotients
of $e^{-t\lambda}$ are dominated by the same type of bound after
slightly enlarging that interval. Dominated convergence in the spectral
measure therefore proves $u'(t)=-HT_tf$ in $L^2$ and continuity of this
derivative. Strong continuity at zero was already established.

The preceding distributional identity, applied to $u(t)$, gives the
tested differential equation. The kernel formula gives a jointly
continuous representative: on a compact set of $(t,x)$ with $t>0$,
its rows have a common Gaussian envelope in $L^2$ as functions of $y$.
Joint kernel continuity and dominated convergence give continuity of
these rows in $L^2$, and their pairing with $f$ preserves continuity.

For uniqueness let $w$ be the difference of two solutions in the
stated class. The lower bound $H\geq-M$ gives, for $0<t<b$,
\[
\frac{d}{dt}\|w(t)\|_2^2
=-2\operatorname{Re}\langle Hw(t),w(t)\rangle
\leq2M\|w(t)\|_2^2.
\]
Consequently $e^{-2Mt}\|w(t)\|_2^2$ is nonincreasing on $(0,b)$.
For $0<s<t<b$ this implies
$\|w(t)\|_2^2\leq e^{2M(t-s)}\|w(s)\|_2^2$.
Let $s\downarrow0$ and use $w(0)=0$ and continuity; then $w(t)=0$.
Continuity also gives equality at $b$. This proves uniqueness without
requiring differentiability at time zero.
\Done

\subsection{First spatial derivatives of the kernel}
Bounded measurable potentials permit more spatial regularity than mere
continuity of the kernel. The following estimate requires no derivative
of the potential. It does not assert the existence of classical second
spatial derivatives.

\begin{proposition}[Continuous first derivatives and Gaussian bounds]
\label{prop:schrodinger_kernel_gradient}
Put $b_d=2^{d/2}/\sqrt e$. For every real $V\in L^\infty$,
the first derivatives of $k_V(t,x,y)$ in $x$ and in $y$ exist and are
jointly continuous for $t>0$. With $M=\|V\|_\infty$,
\[
\begin{split}
|\nabla_x k_V(t,x,y)|,
\ |\nabla_y k_V(t,x,y)|
&\leq B_{d,M}(t)\,p_{2t}(x-y),\\
B_{d,M}(t)
&=b_d e^{Mt}t^{-1/2}\bigl(1+2^{d/2+1}Mt\bigr).
\end{split}
\]
The norm of each gradient is its Euclidean norm; no mixed derivative
is asserted.
\end{proposition}
\Proof
Direct differentiation and comparison of two Gaussian scales give
\[
|\nabla p_r(w)|
=2^{d/2}\frac{|w|}{2r}e^{-|w|^2/(8r)}p_{2r}(w)
\leq b_d r^{-1/2}p_{2r}(w).
\]
Indeed, the maximum of $\rho e^{-\rho^2/(8r)}$ for $\rho\geq0$
is $2\sqrt{r/e}$. Absolute summation of the kernel recursion gives
the pointwise integral equation
\[
k_V(t,x,y)=p_t(x-y)
-\int_0^t\int p_{t-s}(x-u)V(u)k_V(s,u,y)\,du\,ds.
\]
Summation and integration can be interchanged because the sum of the
absolute integrands is bounded after spatial integration by
$M e^{Ms}p_t(x-y)$, which is integrable in $s$.

The candidate derivative integral is absolutely convergent. Gaussian
convolution and $t\leq2t-s\leq2t$ give
\[
\begin{split}
\int |\nabla p_{t-s}(x-u)|\,|V(u)|\,|k_V(s,u,y)|\,du
&\leq b_d M e^{Ms}(t-s)^{-1/2}p_{2t-s}(x-y)\\
&\leq 2^{d/2}b_d M e^{Mt}(t-s)^{-1/2}p_{2t}(x-y).
\end{split}
\]
For the last step, the normalization ratio is at most $2^{d/2}$,
and the Gaussian exponential at time $2t-s$ is at most the one at
time $2t$. Integration in $s$ gives the bound
$2^{d/2+1}b_d M e^{Mt}\sqrt t\,p_{2t}(x-y)$.
Adding the free derivative bound and using $e^{Mt}\geq1$ yields
the asserted coefficient $B_{d,M}(t)$.

We justify both differentiation and continuity, rather than only
estimating a formal derivative. Substitute $s=t\theta$ and initially
restrict to $\varepsilon\leq\theta\leq1-\varepsilon$.
On compact sets of $(t,x,y)$ with $t>0$, the truncated integral and
its derivative in $x$ have integrable Gaussian majorants in $u$.
Their integrands are continuous in $t,x,y$ for each $\theta,u$;
the factor $V(u)$ is fixed. Differentiation under the integral and
dominated convergence therefore give truncated kernels that are
continuously differentiable in $x$, with their $x$-gradients jointly
continuous in $(t,x,y)$.

The omitted parts of the undifferentiated integral tend uniformly
to zero on these compact sets by its bound $M e^{Ms}p_t(x-y)$.
For the derivative, the preceding estimate reduces the endpoint
errors, up to a compact-set constant, to
\[
\int_0^\varepsilon(1-\theta)^{-1/2}\,d\theta
 +\int_{1-\varepsilon}^1(1-\theta)^{-1/2}\,d\theta,
\]
which tends to zero. The truncated gradients thus converge locally
uniformly to a continuous function $G$. To identify $G$ as the
gradient, apply the fundamental theorem of calculus to each truncated
kernel along a coordinate segment. Uniform convergence of both the
functions and their derivatives passes this identity to the limit.
Dividing by the segment length and using continuity of $G$ proves
that $\nabla_x k_V=G$. Finally symmetry
$k_V(t,x,y)=k_V(t,y,x)$ proves the same assertions in $y$.
\Done

Consequently the integral representative of $T_tf$ belongs to
$C^1(\mathbb R^d)$ for every $f\in L^p$, $1\leq p\leq\infty$,
at each fixed $t>0$. On a compact set of $x$, the kernel and its first
derivatives have common Gaussian envelopes in $y$. Their products
with $|f(y)|$ are integrable by H\"older's inequality, including its
endpoints. Differentiation and continuity therefore pass under the
integral. For $1\leq p\leq q\leq\infty$ and
$1+1/q=1/p+1/r$, the convolution estimate already proved gives
\[
\|\nabla T_tf\|_q
\leq B_{d,M}(t)\,\|p_{2t}\|_r\,\|f\|_p.
\]
Here the $L^q$ norm of the gradient uses its pointwise Euclidean norm,
and $\|p_{2t}\|_r$ is the explicit Gaussian norm from
Corollary~\ref{cor:schrodinger_Lp_Lq} with $t$ replaced by $2t$.

\subsection{Uniform initial traces on spaces of continuous functions}
The kernels diverge on the diagonal as time tends to zero, but their
integrals against suitable initial data have a uniform initial trace.
This remains true when the bounded potential is merely measurable.

\begin{proposition}[Strong continuity on $C_0$ and $\operatorname{BUC}$]
\label{prop:schrodinger_continuous_function_semigroups}
Let $V$ be real and essentially bounded, with $M=\|V\|_\infty$.
For $t>0$, define $T_tf$ by the continuous kernel integral, and put
$T_0f=f$. On each of $C_0(\mathbb R^d)$ and
$\operatorname{BUC}(\mathbb R^d)$ these operators form a strongly
continuous semigroup in the supremum norm, with
\[
\|T_t\|\leq e^{Mt},\qquad
\|T_tf-f\|_\infty
\leq\|P_tf-f\|_\infty+(e^{Mt}-1)\|f\|_\infty.
\]
Here $C_0$ consists of continuous functions vanishing at infinity,
and $\operatorname{BUC}$ consists of bounded uniformly continuous
functions. Both carry the supremum norm.
The operator norm in this display is taken on the indicated function space.
\end{proposition}
\Proof
Summing the Dyson estimates for orders $n\geq1$ gives
\[
|k_V(t,x,y)-p_t(x-y)|\leq(e^{Mt}-1)p_t(x-y).
\]
Consequently, for every bounded measurable $f$,
\[
\|T_tf-P_tf\|_\infty\leq(e^{Mt}-1)\|f\|_\infty,
\qquad \|T_tf\|_\infty\leq e^{Mt}\|f\|_\infty.
\]
The gradient bound in Proposition~\ref{prop:schrodinger_kernel_gradient}
and differentiation under the kernel integral give
\[
\sup_x|\nabla T_tf(x)|\leq B_{d,M}(t)\|f\|_\infty.
\]
Integration along line segments therefore makes $T_tf$ globally
Lipschitz for every fixed $t>0$. In particular it belongs to
$\operatorname{BUC}$.

If $f\in C_0$, fix $\epsilon>0$ and choose $R$ so that
$|f(y)|\leq\epsilon$ for $|y|>R$. Gaussian domination gives
\[
|T_tf(x)|\leq e^{Mt}\left[
\epsilon+\|f\|_\infty\int_{|y|\leq R}p_t(x-y)\,dy\right].
\]
For fixed $t,R$, the last integral tends to zero as $|x|\to\infty$:
it is at most the volume of the ball times
$(4\pi t)^{-d/2}\exp(-(|x|-R)^2/(4t))$ when $|x|>R$.
Letting $\epsilon\downarrow0$ proves $T_tf\in C_0$.

For bounded $f$, the double integral defining $T_tT_sf(x)$ is
absolutely convergent: the integral of its absolute integrand is at
most $e^{M(t+s)}\|f\|_\infty$, using Gaussian convolution and mass one.
Fubini and the pointwise kernel composition law thus give
$T_tT_sf=T_{t+s}f$ for $s,t>0$; the law at zero follows by definition.
The earlier free uniform-trace result,
Corollary~\ref{cor:heat_uniform_initial_trace}, and the displayed
perturbation estimate imply $\|T_tf-f\|_\infty\to0$ for
$f\in\operatorname{BUC}$. Every function in $C_0$ is uniformly
continuous: use uniform continuity on a large closed ball, and smallness
of the function outside a slightly smaller ball for pairs outside it.
Thus the same limit holds on $C_0$.
Finally, for $t>s\geq0$, the semigroup law gives
\[
\|T_tf-T_sf\|_\infty
\leq e^{Ms}\|T_{t-s}f-f\|_\infty.
\]
This proves continuity at every time from continuity at zero.
\Done

\begin{corollary}[The maximal space with a uniform initial trace]
\label{cor:schrodinger_maximal_uniform_trace}
For a real bounded measurable potential and $f\in L^\infty$,
\[
\lim_{t\downarrow0}\|T_tf-f\|_\infty=0
\quad\Longleftrightarrow\quad
f\text{ has a bounded uniformly continuous representative}.
\]
In particular, for $f\in C_b$ the uniform initial trace holds exactly
when $f\in\operatorname{BUC}$.
\end{corollary}
\Proof
Sufficiency follows from
Proposition~\ref{prop:schrodinger_continuous_function_semigroups}.
Conversely, suppose the norm limit holds, and let $f_n$ be the
continuous kernel representative of $T_{1/n}f$, for $n\geq1$. The gradient estimate
makes every $f_n$ bounded and uniformly continuous. The assumed limit
makes this sequence Cauchy in $L^\infty$. A bounded continuous function on
$\mathbb R^d$ has supremum norm equal to its essential supremum norm, since
every nonempty open set has positive Lebesgue measure. Applied to
$f_n-f_m$, this shows that the sequence is uniformly Cauchy. Its
uniform limit $g$ is bounded and uniformly continuous, and convergence
in $L^\infty$ gives $g=f$ almost everywhere. If $f$ itself is
continuous, equality almost everywhere between $f$ and $g$ implies
equality everywhere by the same full-support argument.
\Done

Strong continuity on all of $C_b$ or $L^\infty$ would be false already
for $V=0$. The example $f(x)=\sin(x^2)$ in
Section~\ref{sec:free_heat_equation} gives
$\|P_tf-f\|_\infty\geq1$ for every $t>0$ in dimension one.
The difference is continuous, so its supremum equals its essential
supremum; the example applies to both spaces. In higher dimensions use
$f(x)=\sin(x_1^2)$ and the product Gaussian formula.
The proposition concerns the integral semigroup; it requires no claim
that multiplication by a discontinuous $V$ preserves either $C_0$ or
$\operatorname{BUC}$.

\subsection{Time derivatives and the limit of classical spatial regularity}
The time variable has a different regularity mechanism: at positive
times the spectral semigroup can be differentiated between two fixed
smoothing operators. This avoids differentiating the potential.

\begin{proposition}[Continuous time derivatives of every order]
\label{prop:schrodinger_kernel_time_derivatives}
For each integer $n\geq1$, the derivative $\partial_t^n k_V(t,x,y)$
exists and is jointly continuous for $t>0$. It satisfies
\[
|\partial_t^n k_V(t,x,y)|
\leq(2\pi t)^{-d/2}
 \left[e^{Mt}M^n+e^{Mt/2}\left(\frac{2n}{et}\right)^n\right].
\]
This estimate is uniform in $x,y$; it is not an off-diagonal Gaussian
derivative estimate.
\end{proposition}
\Proof
For $s>0$, spectral calculus gives the bounded operator $H^nT_s$ and
\[
\|H^nT_s\|_{2\to2}
\leq M^n e^{Ms}+\left(\frac{n}{es}\right)^n.
\]
On the nonnegative part of the spectrum, maximize
$\lambda^n e^{-s\lambda}$ at $\lambda=n/s$; on $[-M,0)$ use
$|\lambda|^n e^{-s\lambda}\leq M^n e^{Ms}$. On any compact
positive $s$ interval, the same bounds with $n+1$ and $n+2$ control
the scalar Taylor remainders uniformly in $\lambda\geq-M$.
The spectral operator norm estimate therefore proves that $s\mapsto T_s$
is infinitely differentiable in operator norm, with derivative
$(-H)^nT_s$ of order $n$.

Fix $a>0$. The real functions
$r_x(u)=k_V(a,u,x)=k_V(a,x,u)$ belong to $L^2$, depend continuously
on $x$ in that norm, and obey
\[
\|r_x\|_2\leq e^{Ma}\|p_a\|_2,
\qquad \|p_a\|_2^2=(8\pi a)^{-d/2}.
\]
The norm continuity follows from joint kernel continuity and a common
$L^2$ Gaussian envelope when $x$ ranges over a compact set.
The pointwise composition law, applied twice, gives for $t>2a$
\[
k_V(t,x,y)=\langle T_{t-2a}r_y,r_x\rangle.
\]
Here the inner product is linear in its first variable; the reality
and symmetry of the rows explain the displayed order. Since $a$ is
fixed, only the middle operator depends on $t$. Differentiating in
operator norm gives
\[
\partial_t^n k_V(t,x,y)
=\langle(-H)^nT_{t-2a}r_y,r_x\rangle.
\]
Operator-norm continuity and the $L^2$ continuity of the rows prove
joint continuity on $t>2a$. Every positive time has a neighborhood
of this form, proving the assertion throughout $t>0$.

To estimate at a particular time $t$, we may now choose $a=t/4$
in the already proved identity. Cauchy--Schwarz and the operator
bound with $s=t/2$ give
\[
e^{Mt/2}(2\pi t)^{-d/2}
 \left[M^n e^{Mt/2}+\left(\frac{2n}{et}\right)^n\right],
\]
as required. This choice is made after differentiation; the smoothing
rows were held fixed when taking the derivative.
\Done

\begin{example}[A bounded potential need not give continuous second derivatives]
In dimension one, take $V(x)=\mathbf1_{(0,\infty)}(x)$. For every
$t>0$ and $y\in\mathbb R$, the function $x\mapsto k_V(t,x,y)$
is not $C^2$ on any neighborhood of zero, although its first spatial
derivative and all its time derivatives are continuous.
\end{example}
\Proof
We first justify the spatial distributional equation for a kernel
column. Fix $0<a<t$ and write
\[
k_V(t,\cdot,y)=T_{t-a}k_V(a,\cdot,y).
\]
The initial column is in $L^2$, so
Proposition~\ref{prop:schrodinger_heat_solution} applies. Pairing with
a compactly supported smooth test function identifies its strong
$L^2$ derivative with the pointwise time derivative from the preceding
proposition: on the test function's compact support, that derivative
is continuous and uniformly bounded on a compact time interval, so
differentiation passes under the integral. Consequently
\[
\partial_x^2 k_V(t,\cdot,y)
=\partial_t k_V(t,\cdot,y)+V\,k_V(t,\cdot,y)
\]
as distributions.

If the column were $C^2$ near zero, its continuous second derivative
would equal $\partial_t k_V(t,x,y)$ for $x<0$ and
$\partial_t k_V(t,x,y)+k_V(t,x,y)$ for $x>0$. Distributional
equality gives equality almost everywhere, and continuity on each
open half interval makes it pointwise there. The limits of these
two expressions at zero differ by $k_V(t,0,y)$, which is strictly
positive by the Gaussian lower bound. This contradicts continuity
of the second derivative at zero.
\Done

\subsection{The extension to finite measures}
For a fixed $t>0$, the kernel just constructed connects the concrete
semigroup to Theorem~\ref{thm:separate_kernel_criterion}. To make the
domains precise, its integral formula defines a bounded operator
$T_t^{(1)}:L^1(\mathbb R^d)\to L^\infty(\mathbb R^d)$, since
\[
\left|\int k_V(t,x,y)f(y)\,dy\right|
\leq e^{Mt}(4\pi t)^{-d/2}\|f\|_1.
\]
On $L^1\cap L^2$ this agrees with the spectral semigroup $T_t$,
by its proved kernel representation. This intersection is dense in
$L^1$, so $T_t^{(1)}$ is the unique bounded extension of that restriction
with the indicated domain and codomain.

Lebesgue measure has full support, and the kernel is bounded and
jointly, hence separately, continuous. The criterion therefore gives
the unique extension
\[
S_t:\mathcal M(\mathbb R^d)\longrightarrow C_b(\mathbb R^d),
\qquad (S_t\mu)(x)=\int k_V(t,x,y)\,d\mu(y),
\]
which is both bounded in norm and continuous from $\sigma(\mathcal M,C_b)$ to
$\sigma(C_b,\mathcal M)$. In particular
\[
(S_t\delta_y)(x)=k_V(t,x,y),\qquad
\|S_t\|=\|T_t^{(1)}\|_{1\to\infty}
=\sup_{x,y}|k_V(t,x,y)|
\leq e^{Mt}(4\pi t)^{-d/2}.
\]
Thus evaluation on a point mass has a precise meaning even though
the point mass is not an $L^1$ density. Joint continuity still comes
from the Dyson-series proof above; the abstract criterion alone
only yields separate continuity.

\subsection{The Trotter product formula for bounded potentials}
The product formula provides a second way to approximate the semigroup;
the original product-formula paper is \textcite{Trotter1959Product}.
We prove the precise bounded-potential case needed here, rather than
importing a general product theorem with additional domain hypotheses.
In particular the multiplication operator $e^{-hV}$ need not preserve
$D(H_0)$.

\begin{theorem}[Strong product approximation, uniformly on bounded time intervals]
\label{thm:bounded_potential_trotter}
For the operators of this chapter and every $f\in L^2$ and $b>0$,
\[
\lim_{n\to\infty}\sup_{0\leq t\leq b}
\left\|\left(P_{t/n}e^{-(t/n)V}\right)^n f-T_tf\right\|_2=0.
\]
The exponential $e^{-hV}$ is multiplication by the bounded measurable
function $x\mapsto e^{-hV(x)}$.
\end{theorem}
\Proof
Put $F(h)=P_he^{-hV}$ for $h\geq0$. Then $F(0)=I$ and
\[
\|F(h)\|\leq e^{Mh},\qquad \|T_h\|\leq e^{Mh}.
\]
Let $D=D(H_0)=D(H)$, with graph norm
$\|u\|_D=\|u\|_2+\|H_0u\|_2$.
For $u\in D$ split the difference quotient in the order
\[
\frac{F(h)u-u}{h}
 =P_h\frac{e^{-hV}-I}{h}u+\frac{P_hu-u}{h}.
\]
The first term tends to $-Vu$: boundedness of $V$ gives
$(e^{-hV}-I)/h\to-V$ in operator norm, and $P_h$ tends strongly
to $I$. The second tends to $-H_0u$ by spectral calculus on $D(H_0)$.
Likewise $(T_hu-u)/h\to-Hu$. Consequently the linear maps
\[
E_h:D\longrightarrow L^2,\qquad E_hu=\frac{F(h)u-T_hu}{h},
\quad h>0,
\]
converge to zero on each fixed $u\in D$. This calculation never
applies $H_0$ to $e^{-hV}u$.

We need this convergence uniformly on one compact set, so record the
necessary bound. For $0<h\leq b$, scalar spectral estimates and
the semigroup derivative on $D(H)$ give
\[
\begin{split}
\left\|\frac{P_hu-u}{h}\right\|_2&\leq\|H_0u\|_2,\\
\left\|\frac{e^{-hV}-I}{h}u\right\|_2
 &\leq M e^{Mb}\|u\|_2,\\
\left\|\frac{T_hu-u}{h}\right\|_2
 &\leq e^{Mb}\|Hu\|_2
 \leq e^{Mb}\bigl(\|H_0u\|_2+M\|u\|_2\bigr).
\end{split}
\]
For the third displayed estimate, one can integrate
$\frac{d}{ds}T_su=-T_sHu$ from $0$ to $h$; the integrand is
strongly continuous. It follows that
\[
\|E_hu\|_2
\leq(1+e^{Mb})\|H_0u\|_2+2Me^{Mb}\|u\|_2
\leq C_b\|u\|_D
\]
with a constant independent of $h$.

Here is the compactness argument explicitly. If linear maps $E_h$
are uniformly bounded by $C_b$ from $D$ to $L^2$ and converge to
zero at every fixed vector, then they converge uniformly on each
graph-norm compact subset $Q$ of $D$. Indeed, choose a finite
$\varepsilon$-net $u_1,\ldots,u_N$ in $Q$. For $u\in Q$ and a
net point within $\varepsilon$,
\[
\|E_hu\|_2\leq C_b\varepsilon+
                     \max_{1\leq j\leq N}\|E_hu_j\|_2.
\]
The finite maximum tends to zero; then let $\varepsilon$ tend to zero.

Fix $f\in D$ and use
$Q_f=\{T_sf:0\leq s\leq b\}$. This set is compact in the graph
norm. In fact $s\mapsto T_sf$ is continuous in $L^2$, and
\[
H_0T_sf=T_sHf-VT_sf
\]
is continuous there as well, by the spectral domain invariance and
commutation identity $HT_sf=T_sHf$. Thus $Q_f$ is the continuous
image of a compact interval. In particular
\[
\omega_f(\delta)=
\sup_{\substack{u\in Q_f\\0<h\leq\delta}}\|E_hu\|_2
\longrightarrow0\quad(\delta\downarrow0).
\]

For $0<t\leq b$ put $h=t/n$. The algebraic telescoping identity,
which requires no commutation between $F(h)$ and $T_h$, reads
\[
F(h)^nf-T_h^nf
 =\sum_{j=0}^{n-1}F(h)^{n-1-j}(F(h)-T_h)T_h^jf.
\]
Since $T_h^jf=T_{jh}f\in Q_f$ and $T_h^n=T_t$, its norm is at
most
\[
nh\,e^{Mb}\omega_f(b/n)
\leq b e^{Mb}\omega_f(b/n),
\]
which tends to zero independently of $t$. At $t=0$ the difference
vanishes exactly. This proves the theorem on $D$.

Finally $D$ is dense in $L^2$. For $f\in L^2$ and $g\in D$,
the uniform bounds on both approximants give
\[
\sup_{0\leq t\leq b}\|F(t/n)^nf-T_tf\|_2
\leq2e^{Mb}\|f-g\|_2
 +\sup_{0\leq t\leq b}\|F(t/n)^ng-T_tg\|_2.
\]
Take the limit superior in $n$ and then approximate $f$ by such $g$.
\Done

This theorem proves strong operator convergence, uniformly in the
stated time range. It does not assert convergence in operator norm,
nor pointwise convergence of the product kernels. Kernel existence,
joint continuity and Gaussian domination were proved independently
above, so none of them rests on an unjustified interchange of the
product limit with point evaluation.

For comparison with contraction-semigroup formulations, replace $V$ by
$V+M$. Both $P_t$ and multiplication by $e^{-t(V+M)}$ are contractions,
and $H+M\geq0$. The exact scalar identities are
\[
\begin{aligned}
\bigl(P_{t/n}e^{-(t/n)(V+M)}\bigr)^n
 &=e^{-Mt}\bigl(P_{t/n}e^{-(t/n)V}\bigr)^n,\\
e^{-t(H+M)}&=e^{-Mt}T_t.
\end{aligned}
\]
Multiplication by $e^{Mt}$ transfers the shifted approximation to the
original one, preserving uniform strong convergence on each $[0,b]$.

The Gaussian normalization, the factor $e^{t\|V\|_\infty}$, and the
restriction $t>0$ are part of these statements. At time zero the
semigroup is the identity; the heat kernels approach a point mass,
not a continuous function kernel of the same form. Unbounded potentials
require additional hypotheses and are not included in this bounded
perturbation argument.
\ifdefined\qwertyflag
\undef\qwertyflag
\end{document}
\fi
